\Needspace{10\baselineskip}
\section{A Liouville theorem for the limiting cone}
\label{sec:stw-liouville}

Rescaling at a point of large gradient removes the background and first-order terms of the Gauduchon equation and leaves a degenerate equation on $\mathbb C^n$, whose solutions are maximal for a cone of matrices. We show that its bounded Lipschitz solutions are constant. The proof argues by induction on the dimension: if the energy of a solution in some complex direction becomes small on large balls, a limit is independent of one variable; if all directional energies stay positive, they give a strictly admissible comparison function.

On $\mathbb C^n$ we use Euclidean coordinates $z=(z_1,\ldots,z_n)$, the identity matrix $I$, the Euclidean balls $B(x,r)$ and Euclidean volume $\mathrm{d}y$. In this section $\nabla$ and $\Delta$ denote the Euclidean real gradient and Laplacian, and $\partial_\xi$ denotes differentiation in a complex direction $\xi$.

\Needspace{8\baselineskip}
For $n\geq2$ and a Hermitian matrix $H$, set
\[
 P_n(H)=\frac{(\operatorname{tr}H)I-H}{n-1}.
\]
This is the map $P_\gamma$ of \cref{sec:stw-matrix-structure} for the Euclidean metric. It preserves nonnegativity and satisfies $P_n(I)=I$ and $\operatorname{tr}P_n(H)=\operatorname{tr}H$. For a real function $v$ write $H(v)=(v_{i\bar j})$ in Euclidean coordinates. A continuous function $v$ on an open set of $\mathbb C^n$ is \emph{admissible}\index{Liouville theorem!admissible functions} if $P_n(H(v))\geq0$ in the sense of matrix-valued distributions. It is \emph{maximal for this cone} if it is admissible and every smooth function $h$ touching $v$ from below at a point satisfies
\[
 P_n(H(h))\not>0
\]
at that point. Here touching means that $v-h$ has a local minimum with value zero, which need not be a strict minimum. Maximality is thus the viscosity supersolution condition for the degenerate equation $\det P_n(H(v))=0$ on the boundary of the cone, together with distributional admissibility.

Taking the trace shows that admissible functions are subharmonic. Since $P_n$ is linear, convolution preserves admissibility, and so do nonnegative linear combinations and integrals of admissible smooth functions. Translations and unitary changes of coordinates preserve both conditions, as do multiplication by a positive constant and addition of a constant.

\begin{lemma}\label{lem:cone-maximal-stability-comparison}
Local uniform limits of admissible maximal functions are admissible and maximal. If $v$ is maximal on a bounded open set $D\subset\mathbb C^n$, $h$ is smooth and admissible on $D$, both functions are continuous on $\overline D$, and $h\leq v$ on $\partial D$, then $h\leq v$ on $D$.
\end{lemma}
\begin{proof}
Let $v_j$ be admissible maximal functions converging locally uniformly to $v$. Distributional positivity is preserved under this convergence. For maximality, suppose that a smooth function $h$ touches $v$ from below at $z_0$ and $P_n(H(h)(z_0))>0$. Choose a small ball on which $v-h\geq0$. Replacing $h$ by $h-\delta|z-z_0|^2$, with $\delta>0$ sufficiently small, keeps its matrix strictly positive near $z_0$ and makes the minimum strict. Uniform convergence then produces interior minimum points of $v_j-h+\delta|z-z_0|^2$ tending to $z_0$. Adding the minimum value to the test function gives a function touching $v_j$ from below with strictly positive matrix, a contradiction.

For comparison, choose $R$ with $\overline D\subset\{|z|\leq R\}$. If $h>v$ somewhere, then for small $\delta>0$ the smooth function
\[
 h_\delta=h+\delta(|z|^2-R^2)
\]
still exceeds $v$ at some interior point and is at most $v$ on the boundary. Hence the negative minimum of $v-h_\delta$ is attained in the interior. But
$P_n(H(h_\delta))=P_n(H(h))+\delta I>0$ there, which contradicts maximality after the minimum value is added to $h_\delta$. No boundary regularity of $D$ is needed for this argument.
\end{proof}

\begin{theorem}\label{thm:stw-cone-liouville}\index{Liouville theorem!Gauduchon cone}
A bounded globally Lipschitz function on $\mathbb C^n$, $n\geq2$, that is maximal for this cone is constant.
\end{theorem}
\begin{proof}
Let $v$ be a bounded globally Lipschitz maximal function on $\mathbb C^n$. We argue by induction on $n$. For $n=2$, $P_2(H)\geq0$ is equivalent to $H\geq0$, so $v$ is plurisubharmonic, and \cref{lem:psh-liouville} shows that it is constant.

We first explain the reduction in dimension for $n\geq3$. If a function is independent of $z_n$, its complex Hessian has the block form $H\oplus0$. The first block of $P_n(H\oplus0)$ is
\[
 \frac{n-2}{n-1}P_{n-1}(H),
\]
and the last entry is $\operatorname{tr}H/(n-1)$. Nonnegativity of $P_{n-1}(H)$ implies $\operatorname{tr}H\geq0$, and positive definiteness implies $\operatorname{tr}H>0$. Thus $P_n(H\oplus0)$ is nonnegative, respectively positive definite, exactly when $P_{n-1}(H)$ is. The distributional version follows from the same linear identities. A test function in $n-1$ variables extends to $\mathbb C^n$ independently of $z_n$, so maximality also passes to $n-1$ variables. Consequently the induction hypothesis applies to bounded Lipschitz maximal functions that are independent of one complex variable.

Suppose now that $v$ is nonconstant. Adding a constant and multiplying by a positive constant, we may assume that $\inf v=0$ and $\sup v=1$. Let $L$ be the Lipschitz constant of $v$. Choose a smooth radial nonnegative probability kernel $\eta$, supported in the closed unit ball and strictly positive in its interior. For a bounded measurable function $f$ and $r>0$, write
\[
 [f]_r(x)=\int_{\mathbb C^n}f(x+y)\eta_r(y)\,\mathrm{d}y,
 \qquad \eta_r(y)=r^{-2n}\eta(y/r),\qquad v_\epsilon=[v]_\epsilon.
\]
Both $v$ and $v^2$ are subharmonic. For $v^2$, use $v_\epsilon\geq0$ and
$\Delta(v_\epsilon^2)=2v_\epsilon\Delta v_\epsilon+2|\nabla v_\epsilon|^2\geq0$, and pass to the uniform limit as $\epsilon\to0$.

For $f=v$ and for $f=v^2$, the averages $[f]_r(x)$ increase with $r$, by the monotonicity of spherical means of subharmonic functions. Their limit as $r\to\infty$ is independent of $x$, since
\[
 |[f]_r(x)-[f]_r(y)|
 \leq\|f\|_\infty\|\eta_r(\cdot-x)-\eta_r(\cdot-y)\|_{L^1}
 \leq C\frac{|x-y|}{r}.
\]
By the submean inequality this limit is at least $f(x)$ for every $x$, while it is at most $\sup f=1$. Therefore
\begin{equation}\label{eq:cone-large-averages}
 [v]_r(x)\longrightarrow1,\qquad [v^2]_r(x)\longrightarrow1.
\end{equation}

We distinguish two cases. First suppose that there exist $\rho>0$, sequences $\epsilon_k\downarrow0$ and $r_k\to\infty$, points $x_k$ and unit complex directions $\xi_k$ such that
\begin{equation}\label{eq:cone-gap-case}
 [v^2]_{r_k}(x_k)+[v]_\rho(x_k)-2v(x_k)\geq\frac43,
 \qquad
 \int_{B(x_k,r_k)}|\partial_{\xi_k}v_{\epsilon_k}|^2\,\mathrm{d}y
       \longrightarrow0.
\end{equation}
Translate $x_k$ to the origin and make $\xi_k$ the last complex direction by a unitary change of coordinates. The resulting functions have the same bounds and Lipschitz constant, so by the Arzel\`a--Ascoli theorem and \cref{lem:cone-maximal-stability-comparison} a subsequence converges locally uniformly to a maximal function $a$. The corresponding mollifications at scale $\epsilon_k$ have the same local uniform limit, and by the integral condition their derivatives in the last complex direction tend to zero in $L^2$ on every fixed ball. Hence $\partial a/\partial z_n=0$ in the sense of distributions. Since $a$ is real and continuous, it is independent of $z_n$.

The dimension reduction and the induction hypothesis show that $a$ is a constant $c\in[0,1]$. In the first inequality of \eqref{eq:cone-gap-case}, the first term is at most one, while the average at the fixed radius $\rho$ and the value at the origin both tend to $c$. It would follow that $4/3\leq1-c\leq1$, a contradiction.

\Needspace{16\baselineskip}
In the remaining case, for every fixed $\rho>0$ there is $C_\rho\geq1$ such that, whenever $0<\epsilon<C_\rho^{-1}$, $r>C_\rho$, and
\[
 [v^2]_r(x)+[v]_\rho(x)-2v(x)\geq\frac43,
\]
one has
\begin{equation}\label{eq:cone-directional-energy}
 \int_{B(x,r)}|\partial_\xi v_\epsilon|^2\,\mathrm{d}y
       \geq C_\rho^{-1}
       \quad\hbox{for every unit complex direction }\xi.
\end{equation}
Indeed, if this failed for $C_\rho=k$ for every integer $k$, a diagonal choice would give sequences as in \eqref{eq:cone-gap-case}.

After a translation we may assume that $v(0)<1/12$. By \eqref{eq:cone-large-averages}, choose $\rho$ with $[v]_\rho(0)>3/4$, and then $r>C_\rho$ with $[v^2]_r(0)>3/4$. The open set
\[
 D=\{x:2v(x)<[v^2]_r(x)+[v]_\rho(x)-4/3\}
\]
contains the origin. We claim that, for some $c>0$,
\[
 [v^2]_r-c|z|^2
\]
is smooth and admissible on $D$.

To prove the claim, differentiate the square of the smooth mollification:
\[
 H(v_\epsilon^2)
       =2v_\epsilon H(v_\epsilon)
            +2(\partial_i v_\epsilon\,\overline{\partial_jv_\epsilon})_{ij}.
\]
Denote by $A_\epsilon$ and $B_\epsilon$ the convolutions at the fixed radius $r$ of these two terms. The first satisfies $P_n(A_\epsilon)\geq0$, because $v_\epsilon\geq0$, $v_\epsilon$ is admissible, and $P_n$ is linear.

The second satisfies $B_\epsilon\geq c_0I$ on $D$ for small $\epsilon$, with a constant $c_0>0$. Here the passage from the unweighted integral in \eqref{eq:cone-directional-energy} to the convolution needs care. The bound $|\nabla v_\epsilon|\leq L$ makes the energy of a sufficiently thin boundary annulus of $B(x,r)$ at most $C_\rho^{-1}/2$, uniformly in $x,\epsilon,\xi$. On the remaining smaller closed ball, the fixed kernel $\eta_r$ has a positive minimum. For points $x\in D$, \eqref{eq:cone-directional-energy} then gives a uniform positive lower bound for the weighted integral in every direction, which is the asserted matrix inequality for $B_\epsilon$.

Choose $0<c\leq c_0$. Since $P_n$ preserves nonnegative matrices,
\[
 P_n\bigl(H([v_\epsilon^2]_r)-cI\bigr)\geq0
       \quad\hbox{on }D.
\]
As $\epsilon\downarrow0$, $v_\epsilon^2\to v^2$ uniformly, and convolution at the fixed positive radius $r$ gives convergence of all derivatives. This proves the claim.

The smooth function
\[
 h=\tfrac12\bigl([v^2]_r-c|z|^2+[v]_\rho-4/3\bigr)
\]
is therefore admissible on $D$. The set $D'=\{v<h\}$ contains the origin and lies in $D$. It is bounded, because $v\geq0$ while $h\to-\infty$ at infinity, the averages being bounded. Moreover $v=h$ on $\partial D'$. \Cref{lem:cone-maximal-stability-comparison}, applied on $D'$, gives $h\leq v$ there, which contradicts the definition of $D'$. This excludes the second case and proves the theorem.
\end{proof}

The Lipschitz bound is used both for the compactness of translates and for the annulus estimate. For $n>1$, the conclusion would not follow from boundedness and subharmonicity alone, since $\mathbb C^n=\mathbb R^{2n}$ carries nonconstant bounded subharmonic functions.
